\documentclass[10pt,a4paper]{amsart}
\usepackage[margin=19mm]{geometry}
\usepackage{amsmath,amssymb,mathtools}
\usepackage[scr=rsfs,cal=euler]{mathalfa}
\usepackage{xfrac}
\usepackage[unicode,hidelinks,bookmarks=false]{hyperref}
\newcommand{\ie}{i.e.\ }
\newcommand{\cf}{cf.\ }
\newcommand{\resp}{respectively\ }
\newcommand{\Subsection}{\subsection}
\providecommand{\nobreakdash}{\nobreak\hbox{-}}
\setlength{\emergencystretch}{2em}
\multlinegap0pt
\usepackage{amssymb}
\usepackage{bm}
\usepackage{dsfont}
\usepackage{csquotes}
\usepackage{mathtools}
\newtheorem{lemma}{Lemma}
\newcommand{\Eqsref}[2]{Equations~(\ref{#1}) and (\ref{#2})}
\newcommand{\eqsref}[2]{Eqs.~(\ref{#1}) and (\ref{#2})}
\newcommand{\eref}[1]{Eq.~\textup{(\ref{#1})}}
\newcommand{\lref}[1]{Lemma~\ref{#1}}
\newcommand{\mmod}{\!\mod} 
\title{Optimal Shadow Estimation with Minimal Measurement Settings}
\author{Zhiyao Yang, Datong Chen,, Huangjun Zhu}
\date{Source: arXiv:2606.20003v1 (CC BY 4.0)}
\begin{document}
\maketitle

\noindent\emph{Source and licence.} Optimal Shadow Estimation with Minimal Measurement Settings. Version arXiv:2606.20003v1 (CC BY 4.0), distributed by arXiv under the Creative Commons Attribution 4.0 International licence, http://creativecommons.org/licenses/by/4.0/, as recorded in the arXiv licence metadata for this exact version and shown on the abstract page https://arxiv.org/abs/2606.20003v1. Full licence notice: \texttt{licenses/arxiv-2606.20003-CC-BY-4.0.txt}.\par\smallskip

\noindent\textit{Editorial scope.} Counted unit: the article's lemma lem:DesignPermutation (source0.tex lines 865--884). The article's complete original proof of it is delivered with it (source0.tex lines 886--905, one \texttt{proof} environment of 20 lines). No mathematics is written, repaired or re-derived: the statement and the proof are the article's own lines, in the article's own notation, and no retained line is substituted or re-flowed.  No further source-local context is needed: every cross-reference of every retained line resolves inside the two delivered spans.  Nothing was omitted from any retained span, so the package carries no omission record.  No retained line contains a citation, so the wrapper carries no bibliography and no bibliography entry.  No retained line contains a figure environment, an image-inclusion command or drawing code, so no image is delivered and the package declares no asset dependency.  Editorial formatting only, in addition: an AMS \texttt{amsart} preamble that restates the article's own declarations and macros used by the retained lines, namely the environments \texttt{lemma} on separate counters, together with the standard packages those lines need.  The retained environments print in this document as Lemma 1 in order of appearance, whereas the article prints the same items with the numbering of its own full document.  The complete native source of the article as distributed in the arXiv e-print for this exact version is bundled unchanged as \texttt{sources/203184/source0.tex}.
\begin{lemma}\label{lem:DesignPermutation}
Suppose $x_1, x_2, x_3, y_1, y_2, y_3\in\{0,1,2,\ldots, d-1\}$. Then the sequence $y_1, y_2, y_3$ is a permutation of $x_1, x_2, x_3$ if and only if
\begin{equation}\label{eq:ModularCondition}
\begin{aligned}
     x_1+x_2+x_3&=y_1+y_2+y_3\mmod d,\\
f_3(x_1)+f_3(x_2)+f_3(x_3)&=f_3(y_1)+f_3(y_2)+f_3(y_3)\mmod 7,\\
x_1^2+x_2^2+x_3^2&=y_1^2+y_2^2+y_3^2\mmod p,\\
x_1^3+x_2^3+x_3^3&=y_1^3+y_2^3+y_3^3\mmod p.
\end{aligned}
\end{equation}
If $d$ is not divisible by 3, then the sequence $y_1, y_2, y_3$ is a permutation of $x_1, x_2, x_3$ if and only if
\begin{equation}\label{eq:ModularCondition2}
\begin{aligned}
     x_1+x_2+x_3&=y_1+y_2+y_3\mmod d,\\
x_1+x_2+x_3&=y_1+y_2+y_3\mmod 3,\\
x_1^2+x_2^2+x_3^2&=y_1^2+y_2^2+y_3^2\mmod p,\\
x_1^3+x_2^3+x_3^3&=y_1^3+y_2^3+y_3^3\mmod p.
\end{aligned}
\end{equation}
\end{lemma}

\begin{proof}[Proof of \lref{lem:DesignPermutation}]
If the sequence $y_1, y_2, y_3$ is a permutation of $x_1, x_2, x_3$, then \eref{eq:ModularCondition} holds automatically. Conversely, suppose \eref{eq:ModularCondition} holds. Then the second line implies that 
\begin{equation}
f_3(x_1)+f_3(x_2)+f_3(x_3)=f_3(y_1)+f_3(y_2)+f_3(y_3),\quad |y_1+y_2+y_3-x_1-x_2-x_3|<d,
\end{equation}
given that $0\leq f_3(x)\leq 2$ for $x\in \{0,1,2,\ldots, d-1\}$. The first two lines in \eref{eq:ModularCondition} together yield $x_1+x_2+x_3=y_1+y_2+y_3$ and
\begin{equation}\label{eq:ElementaryPoly1}
x_1+x_2+x_3=y_1+y_2+y_3\mmod p.
\end{equation}
In conjunction with the last two lines in \eref{eq:ModularCondition}, we can deduce that
\begin{equation}\label{eq:ElementaryPoly23}
\begin{aligned}
x_1 x_2 +x_2 x_3+x_3x_1 &= y_1 y_2 +y_2 y_3+y_3y_1 \mmod p, \\
x_1 x_2 x_3 &= y_1 y_2 y_3 \mmod p.
\end{aligned}
\end{equation}
\Eqsref{eq:ElementaryPoly1}{eq:ElementaryPoly23} together imply that the sequence $y_1, y_2, y_3$ is a permutation of $x_1, x_2, x_3$.

Next, suppose $d$ is not divisible by 3; then $d$ and 3 are coprime. If the sequence $y_1, y_2, y_3$ is a permutation of $x_1, x_2, x_3$, then \eref{eq:ModularCondition2} holds automatically. Conversely, suppose \eref{eq:ModularCondition2} holds. The first two lines together give $x_1+x_2+x_3=y_1+y_2+y_3 \mmod 3d$, which in turn implies that $x_1+x_2+x_3=y_1+y_2+y_3$. So \eqsref{eq:ElementaryPoly1}{eq:ElementaryPoly23} hold as before, and the sequence $y_1, y_2, y_3$ is a permutation of $x_1, x_2, x_3$.
\end{proof}

\end{document}
